\setcounter{section}{11}

  \zwischenueberschrift{Soal latihan}

\inputaufgabe
{}
{

Tunjukkan bahwa
\definitionsverweis {produk}{}{}

\mathl{M \times N}{} dari dua
\definitionsverweis {manifold diferensiabel}{}{}
\mathkor {} {M} {dan} {N} {}
sendiri merupakan sebuah manifold diferensiabel.

}
{} {}

\inputaufgabe
{}
{

Misalkan
\mathkor {} {M_1 \subseteq N_1} {dan} {M_2 \subseteq N_2} {}
\definitionsverweis {submanifold tertutup}{}{.} Tunjukkan bahwa
\definitionsverweis {produk}{}{}
\mathl{M_1 \times M_2}{} merupakan sebuah submanifold tertutup dari
\mathl{N_1 \times N_2}{}.

}
{} {}

\inputaufgabe
{}
{

Deskripsikan peta-peta pada
\definitionsverweis {torus}{}{}

\mathl{S^1 \times S^1}{,} yang diperoleh dari proyeksi-proyeksi stereografik.

}
{} {}

\inputaufgabe
{}
{

Tunjukkan bahwa
\mathl{\R^2 \setminus \{0\}}{}
\definitionsverweis {difeomorfik}{}{}
dengan suatu produk dari
\definitionsverweis {manifold-manifold}{}{}
berdimensi satu.

}
{} {}

\inputaufgabe
{}
{

Tunjukkan bahwa
\definitionsverweis {produk}{}{}

\mathl{M \times N}{} dari dua
\definitionsverweis {manifold yang terhubung lintasan}{}{} dan
\definitionsverweis {manifold diferensiabel}{}{}
\mathkor {} {M} {dan} {N} {}
juga terhubung lintasan.

}
{} {}

\inputaufgabe
{}
{

Misalkan $M$ sebuah
\definitionsverweis {manifold diferensiabel}{}{}
dan
\maabbeledisp {\varphi} { M } { M \times M
} { x } {(x,x)
} {,}
merupakan
\definitionsverweis {pemetaan diagonal}{}{}
ke dalam
\definitionsverweis {produk}{}{}

\mathl{M \times M}{.} Tunjukkan bahwa
\definitionsverweis {diagonal}{}{}
$\varphi(M)$ merupakan sebuah
\definitionsverweis {submanifold tertutup}{}{}.

}
{} {}

\inputaufgabe
{}
{

Misalkan $M$ sebuah
\definitionsverweis {manifold diferensiabel}{}{.}
Tunjukkan bahwa pemetaan pertukaran
\maabbeledisp {} {M \times M} { M \times M
} {(P,Q)} {(Q,P)
} {,}
merupakan sebuah
\definitionsverweis {difeomorfisme}{}{}.

}
{} {}

\inputaufgabe
{}
{

Deskripsikan torus
\mathl{S^1 \times S^1}{} sebagai
\definitionsverweis {himpunan hasil rotasi}{}{}
di dalam $\R^3$.

}
{} {}

\inputaufgabe
{}
{

Misalkan

\mavergleichskette
{\vergleichskette
{ R
}
{ > }{ 0
}
{  }{
}
{    }{
}
{   }{
}
}
{}{}{}
dan perhatikan pemetaan
\maabbeledisp {} { \R^3} { \R
} { (x,y,z) } { \left(  \sqrt{x^2+y^2}- R  \right)^2 +z^2
} {.}
Pertama, tunjukkan bahwa pemetaan yang ditampilkan tidak diferensiabel pada seluruh sumbu $z$. Pada himpunan terbuka
$U=\{(x,y,z)\in\R^3\mid x^2+y^2>0\}$, tentukan
\definitionsverweis {titik-titik reguler}{}{}
dan lokus kritis dari pembatasannya, serta bentuk serat di atas

\mavergleichskette
{\vergleichskette
{ s
}
{ \in }{ \R
}
{  }{
}
{    }{
}
{   }{
}
}
{}{}{.}
Nyatakan pula apa yang terjadi untuk $s<0$. Bagaimana bentuk itu berubah ketika $s$ melewati $R^2$, yakni ketika beralih dari

\mavergleichskette
{\vergleichskette
{ \sqrt{s}
}
{ < }{ R
}
{  }{
}
{    }{
}
{   }{
}
}
{}{}{}
ke

\mavergleichskette
{\vergleichskette
{ \sqrt{s}
}
{ > }{ R
}
{  }{
}
{    }{
}
{   }{
}
}
{}{}{?}

}
{} {}

\inputaufgabegibtloesung
{}
{

Buatlah sebuah animasi yang mengilustrasikan
[[Torus/Faserrealisierung/Andere Fasern/Aufgabe|Soal 11.9]].

}
{} {}

\inputaufgabe
{}
{

Definisikan pemetaan
\maabbdisp {} {S^1 \times S^1} { S^2
} {,}
yang, bagi suatu pasangan sudut
\mathl{(\alpha, \beta)}{}, menafsirkan komponen pertama sebagai titik pada khatulistiwa, lalu bergerak dari titik itu sepanjang lingkaran besar menuju utara sesuai dengan komponen kedua. Apakah pemetaan tersebut
\definitionsverweis {diferensiabel}{}{?}
Bagaimanakah bentuk serat-serat pemetaan itu?

}
{} {}

\inputaufgabe
{}
{

Berikan sebuah argumen heuristik bahwa permukaan bola satuan
\mathl{S^2}{} dan torus
\mathl{S^1 \times S^1}{} tidak
\definitionsverweis {homeomorfik}{}{}.

}
{} {}

\inputaufgabe
{}
{

Dengan
\definitionsverweis {manifold diferensiabel}{}{}
manakah
\mathl{S^1 \times S^1 \setminus \triangle}{,} yaitu torus tanpa
\definitionsverweis {diagonal}{}{,}
\definitionsverweis {difeomorfik}{}{?}

}
{} {}

\inputaufgabegibtloesung
{}
{

Misalkan $X$ sebuah
\definitionsverweis {torus}{}{.}
Berikan sebuah
\definitionsverweis {pemetaan diferensiabel}{}{}
surjektif
\maabbdisp {\varphi} {\R^2} {X
} {}
sedemikian sehingga
\definitionsverweis {pemetaan singgung}{}{}
\maabbdisp {T_P(\varphi)} { T_P\R^2 } { T_{\varphi(P)}X
} {}
juga surjektif di setiap titik

\mavergleichskette
{\vergleichskette
{ P
}
{ \in }{ \R^2
}
{  }{
}
{    }{
}
{   }{
}
}
{}{}{}.

}
{} {}

Dalam soal-soal berikut diperkenalkan produk relatif dari himpunan-himpunan dan ruang-ruang topologis.

\inputaufgabe
{}
{

Misalkan
\mathl{L_1,L_2,M}{} himpunan-himpunan dan
\maabbdisp {p_1} {L_1 } { M
} {}
serta
\maabbdisp {p_2} {L_2 } { M
} {}
pemetaan-pemetaan. Kita definisikan

\mavergleichskettedisp
{\vergleichskette
{L_1 \times_M L_2
}
{ \defeq} { \left\{ (x_1,x_2)  \mid   p_1(x_1) = p_2(x_2)         \right\}
}
{ \subseteq} { L_1 \times L_2
}
{ } {
}
{ } {
}
}
{}{}{.}
\aufzaehlungzwei {Tunjukkan bahwa terdapat sebuah diagram komutatif

\mathdisp {\begin{matrix} L_1 \times_M L_2  & \stackrel{  }{\longrightarrow} &  L_1  &  \\ \downarrow & & \downarrow  & \\  L_2  & \stackrel{  }{\longrightarrow} &  M & \! \! \! \!\!\!\!\!   \\ \end{matrix}} {  }

.
} {Misalkan $T$ sebuah himpunan lain dan
\maabb {\psi_1} { T } { L_1
} {}
serta
\maabb {\psi_2} { T } { L_2
} {}
pemetaan-pemetaan yang memenuhi

\mavergleichskettedisp
{\vergleichskette
{ p_1 \circ \psi_1
}
{ =} { p_2 \circ \psi_2
}
{ } {
}
{ } {
}
{ } {}
}
{}{}{.}
Tunjukkan bahwa terdapat tepat satu pemetaan
\maabbdisp {\psi} { T } { L_1 \times_M L_2
} {}
sedemikian sehingga proyeksinya ke
\mathkor {} {L_1} {dan} {L_2} {}
masing-masing berimpit dengan $\psi_1,\psi_2$.
}

}
{} {}

\inputaufgabe
{}
{

Misalkan
\mathl{L_1,L_2,M}{}
\definitionsverweis {ruang-ruang topologis}{}{}
dan
\maabbdisp {p_1} {L_1 } { M
} {}
serta
\maabbdisp {p_2} {L_2 } { M
} {}
\definitionsverweis {pemetaan-pemetaan kontinu}{}{.}
Kita definisikan

\mavergleichskettedisp
{\vergleichskette
{L_1 \times_M L_2
}
{ \defeq} { \left\{ (x_1,x_2)  \mid   p_1(x_1) = p_2(x_2)         \right\}
}
{ \subseteq} { L_1 \times L_2
}
{ } {
}
{ } {
}
}
{}{}{}
dengan
\definitionsverweis {topologi yang diinduksi dari ruang produk di atas}{}{.}
\aufzaehlungzwei {Tunjukkan bahwa terdapat sebuah diagram komutatif

\mathdisp {\begin{matrix} L_1 \times_M L_2  & \stackrel{  }{\longrightarrow} &  L_1  &  \\ \downarrow & & \downarrow  & \\  L_2  & \stackrel{  }{\longrightarrow} &  M & \! \! \! \!\!\!\!\!   \\ \end{matrix}} {  }

dengan pemetaan-pemetaan kontinu.
} {Misalkan $T$ sebuah ruang topologis lain dan
\maabb {\psi_1} { T } { L_1
} {}
serta
\maabb {\psi_2} { T } { L_2
} {}
pemetaan-pemetaan kontinu yang memenuhi

\mavergleichskettedisp
{\vergleichskette
{ p_1 \circ \psi_1
}
{ =} { p_2 \circ \psi_2
}
{ } {
}
{ } {
}
{ } {}
}
{}{}{.}
Tunjukkan bahwa terdapat tepat satu pemetaan kontinu
\maabbdisp {\psi} { T } { L_1 \times_M L_2
} {}
sedemikian sehingga proyeksinya ke
\mathkor {} {L_1} {dan} {L_2} {}
masing-masing berimpit dengan $\psi_1,\psi_2$.
}

}
{} {}

Dalam situasi kedua soal sebelumnya, $L_1 \times_M L_2$
\zusatzklammer {dengan proyeksi ke $L_2$} {} {}
juga dilambangkan dengan $p_2^*L_1$ atau dengan
\mathl{p_1^*L_2}{.} Gagasannya ialah bahwa objek
\maabb {p_1} {L_1} {M
} {}
ditarik balik sepanjang $p_2$ ke basis baru $L_2$.

\inputaufgabe
{}
{

Misalkan
\maabb {p} { Y } { X
} {}
sebuah
\definitionsverweis {pemetaan kontinu}{}{}
di antara
\definitionsverweis {ruang-ruang topologis}{}{}
\mathkor {} {X} {dan} {Y} {.}
Tunjukkan bahwa untuk
\definitionsverweis {produk relatif}{}{}
berlaku relasi

\mavergleichskettedisp
{\vergleichskette
{ X \times_X Y
}
{ \cong} { Y
}
{ } {
}
{ } {
}
{ } {
}
}
{}{}{}
\zusatzklammer {dengan
\maabb {} { X } { X
} {}
menyatakan identitas} {} {.}

}
{} {}

\inputaufgabe
{}
{

Misalkan
\maabb {p} { Y } { X
} {}
sebuah
\definitionsverweis {pemetaan kontinu}{}{}
di antara
\definitionsverweis {ruang-ruang topologis}{}{}
\mathkor {} {X} {dan} {Y} {.}
Misalkan

\mavergleichskette
{\vergleichskette
{ x
}
{ \in }{ X
}
{  }{
}
{    }{
}
{   }{
}
}
{}{}{}
sebuah titik dengan pemetaan yang bersesuaian
\maabbdisp {\iota} {\{x\}} { X
} {.}
Tunjukkan bahwa untuk
\definitionsverweis {produk relatif}{}{}
berlaku relasi

\mavergleichskettedisp
{\vergleichskette
{ \{x\}  \times_X Y
}
{ \cong} { Y_x
}
{ } {
}
{ } {
}
{ } {
}
}
{}{}{,}
yaitu produk relatif tersebut berimpit dengan
\definitionsverweis {serat}{}{}.

}
{} {}

\inputaufgabe
{}
{

Misalkan
\mathkor {} {X,U} {dan} {V} {}
\definitionsverweis {ruang-ruang topologis}{}{}
dan misalkan
\maabb {p_1} {X \times U} {X
} {}
serta
\maabb {p_2} {X \times V} {X
} {}
adalah proyeksi kanonik ke $X$. Tunjukkan bahwa untuk
\definitionsverweis {produk relatif}{}{}
berlaku relasi

\mavergleichskettedisp
{\vergleichskette
{ \left(  X \times U  \right) \times_X \left(  X \times V  \right)
}
{ \cong} { X \times U \times V
}
{ } {
}
{ } {
}
{ } {
}
}
{}{}{.}

}
{} {}

\inputaufgabe
{}
{

Misalkan
\maabb {f,g} {\R} {\R
} {}
\definitionsverweis {fungsi-fungsi kontinu}{}{.}
\aufzaehlungdrei{Tunjukkan bahwa untuk
\definitionsverweis {produk relatif}{}{}
\zusatzklammer {fungsi-fungsi tersebut tidak muncul secara eksplisit dalam notasi} {} {}

\mavergleichskettedisp
{\vergleichskette
{ \R \times_\R \R
}
{ =} { \left\{ (x,y) \in \R^2   \mid   f(x) = g(y)         \right\}
}
{ } {
}
{ } {
}
{ } {
}
}
{}{}{}
berlaku.
}{Sketsalah
\mathl{\R \times_\R \R}{} untuk beberapa pilihan fungsi
\mathl{(f,g)}{.}
}{Tentukan pada (2) apakah
\mathl{\R \times_\R \R}{}
memiliki
\definitionsverweis {titik-titik singular}{}{}.
}

}
{} {}

\inputaufgabe
{}
{

Misalkan
\mathkor {} {X,Y} {dan} {Z} {}
\definitionsverweis {ruang-ruang topologis}{}{}
dan misalkan
\maabb {\varphi} { Y } { X
} {}
serta
\maabb {p} { Z } { X
} {}
\definitionsverweis {pemetaan-pemetaan kontinu}{}{.}
Misalkan
\maabbdisp {p_Y} {Y \times_XZ } { Y
} {.}
Tunjukkan bahwa sebuah
\definitionsverweis {penampang kontinu}{}{}
\maabb {s} { Y } { Y \times_XZ
} {}
setara dengan sebuah pemetaan kontinu
\maabb {t} { Y } { Z
} {}
yang memenuhi

\mavergleichskette
{\vergleichskette
{ p \circ t
}
{ = }{ \varphi
}
{  }{
}
{    }{
}
{   }{
}
}
{}{}{.}

}
{} {}

\inputaufgabe
{}
{

Misalkan
\mathl{X,Y,Z}{} dan $X'$
\definitionsverweis {ruang-ruang topologis}{}{}
dan diberikan sebuah diagram komutatif dari
\definitionsverweis {pemetaan-pemetaan kontinu}{}{}

\mathdisp {\begin{matrix}  & Y& \stackrel{}{\longrightarrow} &Z  \\  &  & \searrow& \downarrow   \\  & X' & \stackrel{\varphi}{  \longrightarrow } & X    \end{matrix}} {  }

. Misalkan

\mavergleichskette
{\vergleichskette
{ Y'
}
{ = }{ \varphi^*Y
}
{ = }{ Y \times_X X'
}
{    }{
}
{   }{
}
}
{}{}{}
dan

\mavergleichskette
{\vergleichskette
{ Z'
}
{ = }{ \varphi^*Z
}
{ = }{ Z \times_X X'
}
{    }{
}
{   }{
}
}
{}{}{.}
Tunjukkan bahwa terdapat diagram komutatif

\mathdisp {\begin{matrix}  & Y' & \stackrel{}{\longrightarrow} & Z'  \\  &  & \searrow& \downarrow   \\  &  &  & X'    \end{matrix}} {  }

.

}
{} {}

\inputaufgabe
{}
{

Misalkan
\mathkor {} {X,Y,Z} {dan} {W} {}
\definitionsverweis {ruang-ruang topologis}{}{}
dan misalkan
\maabb {} { W } { X
} {}
serta

\mathdisp {Z \stackrel{\varphi}{\longrightarrow} Y\stackrel{\psi}{\longrightarrow} X} {  }

merupakan
\definitionsverweis {pemetaan-pemetaan kontinu}{}{.}
Tunjukkan bahwa
\definitionsverweis {produk relatif}{}{}
\zusatzklammer {atau, secara ekuivalen, tarik balik} {} {}
memenuhi isomorfisme-isomorfisme kanonik

\mavergleichskettedisp
{\vergleichskette
{ \varphi^* \left(  \psi^*W  \right)
}
{ \cong} { Z \times_Y \left(  Y \times_X W  \right)
}
{ \cong} { Z \times_XW
}
{ \cong} { ( \psi \circ \varphi)^*W
}
{ } {
}
}
{}{}{,}
yang pada representasi produk diberikan oleh
\mathl{(z,(\varphi(z),w))\longmapsto(z,w)}{.}

}
{} {}

\inputaufgabe
{}
{

Perhatikan
\definitionsverweis {lingkaran}{}{}
$S^1$. Definisikan sebuah \stichwort {struktur grup diferensiabel} {} pada $S^1$, yaitu sebuah unsur identitas
\mathl{P \in S^1}{,} sebuah
\definitionsverweis {pemetaan diferensiabel}{}{}
\maabbeledisp {\alpha} {S^1} {S^1
} { x } { \alpha (x)
} {,}
dan sebuah pemetaan diferensiabel
\maabbeledisp {} {T = S^1 \times S^1} {S^1
} {(x,y)} { \varphi(x,y)
} {,}
sedemikian sehingga $S^1$ menjadi sebuah
\definitionsverweis {grup komutatif}{}{}
dengan data tersebut.

}
{} {}

\inputaufgabe
{}
{

Perhatikan
\definitionsverweis {grup linear umum}{}{}

\mavergleichskette
{\vergleichskette
{G
}
{ = }{ \operatorname{GL}_{  n   } \! \left(  \R  \right)
}
{ \subseteq }{ \R^{n^2}
}
{    }{
}
{   }{
}
}
{}{}{}
sebagai
\definitionsverweis {submanifold terbuka}{}{}
dari $\R^{n^2}$. Definisikan sebuah \stichwort {struktur grup diferensiabel} {} pada $G$, yaitu sebuah unsur identitas
\mathl{P \in G}{,} sebuah
\definitionsverweis {pemetaan diferensiabel}{}{}
\maabbeledisp {\alpha} { G } { G
} { x } { \alpha (x)
} {,}
dan sebuah pemetaan diferensiabel
\maabbeledisp {\varphi} {G \times G} {G
} {(x,y)} { \varphi(x,y)
} {,}
sedemikian sehingga $G$ menjadi sebuah
\definitionsverweis {grup}{}{}
dengan data tersebut.

}
{} {}

\inputaufgabe
{}
{

Tunjukkan bahwa
\definitionsverweis {pemetaan}{}{}
\maabbeledisp {\varphi} {S^1 \times \R } { S^1 \times \R
} {(x_1 , x_2;t)} {(-x_1 , -x_2; -t)
} {,}
merupakan sebuah
\definitionsverweis {difeomorfisme}{}{}
\definitionsverweis {tanpa titik tetap}{}{} yang inversnya adalah dirinya sendiri.

}
{} {}

\inputaufgabe
{}
{

Misalkan

\mavergleichskette
{\vergleichskette
{ D
}
{ = }{ x^2 \partial_x +  \left(  x-y^2  \right) \partial_y
}
{  }{
}
{    }{
}
{   }{
}
}
{}{}{.}
Hitung

\mathdisp {D  \left(  4x^2-xy+5y^3  \right)} { . }

}
{} {}

\inputaufgabe
{}
{

Misalkan

\mavergleichskette
{\vergleichskette
{ D
}
{ = }{ \sum_{j = 1}^n g_j \partial_j
}
{  }{
}
{    }{
}
{   }{
}
}
{}{}{}
sebuah operator diferensial pada suatu himpunan terbuka

\mavergleichskette
{\vergleichskette
{ U
}
{ \subseteq }{ \R^n
}
{  }{
}
{    }{
}
{   }{
}
}
{}{}{.}
Penerapan $D$ pada fungsi-fungsi diferensiabel mana yang menghasilkan fungsi-fungsi koefisien $g_j$?

}
{} {}

\inputaufgabe
{}
{

Untuk suatu himpunan terbuka $U\subseteq\R^n$, misalkan fungsi-fungsi $g_j\in C^0(U)$ diberikan, dan definisikan

\mavergleichskette
{\vergleichskette
{ D
}
{ = }{ \sum_{j = 1}^n g_j \partial_j
}
{  }{
}
{    }{
}
{   }{
}
}
{}{}{}
dengan

\mavergleichskette
{\vergleichskette
{ U
}
{ \subseteq }{ \R^n
}
{  }{
}
{    }{
}
{   }{
}
}
{}{}{.}
Pada setiap fungsi diferensiabel, $D$ bekerja menurut rumus di atas. Tunjukkan bahwa $D$ merupakan sebuah operator diferensial orde pertama dalam pengertian
[[Mannigfaltigkeit/Differentialoperator/Erste Ordnung/Definition|Definisi 11.4]].

}
{} {}

\inputaufgabe
{}
{

Misalkan $M$ sebuah
\definitionsverweis {manifold diferensiabel}{}{}
dan $D$ sebuah
\definitionsverweis {medan vektor}{}{}
diferensiabel pada $M$. Tunjukkan bahwa

\mavergleichskette
{\vergleichskette
{ [ D,D]
}
{ =} { 0
}
{  }{
}
{    }{
}
{   }{
}
}
{}{}{.}

}
{} {}

\inputaufgabe
{}
{

Misalkan $M$ sebuah
$C^3$-\definitionsverweis {manifold diferensiabel}{}{.}
Tunjukkan bahwa untuk
$C^2$-\definitionsverweis {medan-medan vektor}{}{}
$D,E,F$ pada $M$ berlaku apa yang disebut \stichwort {identitas Jacobi} {},

\mavergleichskettedisp
{\vergleichskette
{ [ [D,E],F] + [ [E,F],D] + [ [F,D],E]
}
{ =} { 0
}
{ } {
}
{ } {
}
{ } {
}
}
{}{}{.}

}
{} {}

Misalkan $V$ sebuah
\definitionsverweis {ruang vektor}{}{}
atas suatu
\definitionsverweis {lapangan}{}{}
$K$, dan misalkan
\maabbeledisp {} {V \times V} { V
} {(f,g)} { [f,g]
} {,}
sebuah
\definitionsverweis {operasi biner}{}{}
pada $V$. Pasangan
\mathl{(V, [-,-])}{} disebut sebuah
\definitionswort {aljabar Lie}{,}
jika ketiga syarat berikut dipenuhi.
\aufzaehlungdrei{ $[-,-]$ bersifat
\definitionsverweis {bilinear}{}{.}
}{ $[-,-]$ memenuhi apa yang disebut \stichwort {identitas Jacobi} {,} yaitu untuk

\mavergleichskette
{\vergleichskette
{ u,v,w
}
{ \in }{ V
}
{  }{
}
{    }{
}
{   }{
}
}
{}{}{}
berlaku

\mavergleichskettedisp
{\vergleichskette
{ [ [u,v],w] + [ [v,w],u] + [ [w,u],v]
}
{ =} { 0
}
{ } {
}
{ } {
}
{ } {
}
}
{}{}{.}
}{Berlaku

\mavergleichskette
{\vergleichskette
{ [v,v]
}
{ = }{ 0
}
{  }{
}
{    }{
}
{   }{
}
}
{}{}{}
untuk setiap

\mavergleichskette
{\vergleichskette
{ v
}
{ \in }{ V
}
{  }{
}
{    }{
}
{   }{
}
}
{}{}{.}
}

\inputaufgabe
{}
{

Misalkan $M$ sebuah
$C^\infty$-\definitionsverweis {manifold}{}{}
dan misalkan $V$ ruang semua
$C^\infty$-\definitionsverweis {medan vektor}{}{}
pada $M$. Tunjukkan bahwa $V$, dengan
\definitionsverweis {braket Lie}{}{}
medan-medan vektor, merupakan sebuah
\definitionsverweis {aljabar Lie}{}{}.

}
{} {}

\inputaufgabe
{}
{

Misalkan $V$ sebuah
$K$-\definitionsverweis {ruang vektor}{}{}
atas suatu
\definitionsverweis {lapangan}{}{}
$K$ dan misalkan

\mavergleichskettedisp
{\vergleichskette
{ L
}
{ =} { \operatorname{End}_{  } \, ( V )
}
{ } {
}
{ } {
}
{ } {
}
}
{}{}{}
adalah himpunan semua
$K$-\definitionsverweis {endomorfisme}{}{}
linear dari $V$ ke $V$. Tunjukkan bahwa $L$, yang dilengkapi dengan

\mavergleichskettedisp
{\vergleichskette
{ [f,g]
}
{ \defeq} { f \circ g -g \circ f
}
{ } {
}
{ } {
}
{ } {
}
}
{}{}{,}
merupakan sebuah
\definitionsverweis {aljabar Lie}{}{}.

}
{} {}

  \zwischenueberschrift{Soal untuk dikumpulkan}

\inputaufgabe
{5}
{

Misalkan

\mavergleichskette
{\vergleichskette
{ 0
}
{ < }{ r
}
{ < }{ R
}
{    }{
}
{   }{
}
}
{}{}{}
Identifikasikan $S^1 \cong \R/(2\pi\Z)$ melalui kelas-kelas sudut. Dengan menuliskan unsur-unsurnya sebagai $[\varphi]$ dan $[\psi]$, misalkan

\mavergleichskettedisp
{\vergleichskette
{ T
}
{ =} { \left\{ (x,y,z) \in \R^3  \mid   \left(  \sqrt{x^2+y^2} -R  \right)^2 +z^2 = r^2         \right\}
}
{ } {
}
{ } {
}
{ } {
}
}
{}{}{.}
Tunjukkan bahwa pemetaan
\maabbeledisp {} { S^1 \times S^1 } { T
} { ( [\varphi], [\psi]) } {( (R +r \cos  \psi) \cos \varphi, (R+ r \cos  \psi ) \sin \varphi ,  r \sin  \psi )
} {}
merupakan sebuah
\definitionsverweis {bijeksi}{}{}.

}
{} {}

\inputaufgabe
{6}
{

Misalkan $T$ sebuah
\definitionsverweis {torus}{}{}
dan misalkan
\mathl{P,Q \in T}{} dua titik. Tunjukkan bahwa terdapat suatu lingkungan koordinat bersama
\mathl{P,Q \in U \subseteq T}{} sedemikian sehingga pemetaan koordinat
\maabbdisp {\alpha} { U } { V
} {}
merupakan sebuah
\definitionsverweis {homeomorfisme}{}{}
dengan
\mathl{V= {]0,1[} \times  {]0,1[}}{}.

}
{} {}

\inputaufgabe
{3}
{

Tunjukkan bahwa
\definitionsverweis {produk relatif}{}{}
$\R \times_\R \R$ untuk fungsi-fungsi yang terdiferensialkan secara kontinu,
\maabbdisp {f,g} {\R} {\R
} {}
tidak harus merupakan sebuah
\definitionsverweis {manifold topologis}{}{}.

}
{} {}

\inputaufgabe
{2}
{

Misalkan

\mavergleichskettedisp
{\vergleichskette
{ D
}
{ =} { \left(  xy^2-  \sin z  \right)  \partial_x +  \left(  e^{x+y+z}  \right) \partial_y+  \left(  xy^3z  \right) \partial_z
}
{ } {
}
{ } {
}
{ } {
}
}
{}{}{.}
Hitung

\mathdisp {D  \left(  7e^{xy}-  \cos \left( z^3+y^5 \right) -x^4y^3 + y^2z^3e^{z}  \right)} { . }

}
{} {}

\inputaufgabe
{2}
{

Misalkan

\mavergleichskette
{\vergleichskette
{ W
}
{ \subseteq }{ \R^n
}
{  }{
}
{    }{
}
{   }{
}
}
{}{}{}
\definitionsverweis {terbuka}{}{,}
\maabb {h} { W } { \R
} {}
sebuah
\definitionsverweis {fungsi yang terdiferensialkan secara kontinu}{}{}
dan

\mavergleichskette
{\vergleichskette
{ Y
}
{ = }{ h^{-1}(0)
}
{  }{
}
{    }{
}
{   }{
}
}
{}{}{}
adalah
\definitionsverweis {serat}{}{}
di atas $0$, dengan $h$
\definitionsverweis {reguler}{}{}
di setiap titik $Y$ dan karena itu mendefinisikan sebuah manifold diferensiabel. Misalkan $D$ sebuah operator diferensial orde pertama yang bersesuaian dengan suatu medan vektor kontinu pada $Y$. Tunjukkan bahwa

\mavergleichskettedisp
{\vergleichskette
{ D(gh)
}
{ =} { 0
}
{ } {
}
{ } {
}
{ } {
}
}
{}{}{}
untuk setiap fungsi diferensiabel $g$ pada $Y$.

}
{} {}

\inputaufgabe
{4 (1+1+1+1)}
{

Pada $\R^3$, perhatikan
\definitionsverweis {medan-medan vektor}{}{}
berikut,

\mavergleichskettedisp
{\vergleichskette
{ D
}
{ =} { -y \partial_x  + z\partial_y -x \partial_z
}
{ } {
}
{ } {
}
{ } {
}
}
{}{}{,}

\mavergleichskettedisp
{\vergleichskette
{ E
}
{ =} { 5 \partial_x - 3\partial_y +  \cos \left(  xy \right)  \partial_z
}
{ } {
}
{ } {
}
{ } {
}
}
{}{}{,}
dan

\mavergleichskettedisp
{\vergleichskette
{ F
}
{ =} { 4e^{xyz} \partial_y -z^3 \partial_z
}
{ } {
}
{ } {
}
{ } {
}
}
{}{}{.}
\aufzaehlungvier{Hitung
\mathl{[D,E]}{.}
}{Hitung
\mathl{[D,F]}{.}
}{Hitung
\mathl{[E,F]}{.}
}{Hitung
\mathl{[ [D,E], F]}{.}
}

}
{} {}

 [[Kategorie:Latexseite]]
